\begin{remark}{4.7.1.2}\label{I.4.7.1.2}
--- {\renewcommand{\thefootnote}{(39)}\nde{This remark has been added.}}
Let \(\mathcal{F}\) be a \(G\)-\(\OS\)-module. The subsheaf of invariants
\(\mathcal{F}^{G}\) is defined as follows: for every open \(U\) of \(S\),
\[
\mathcal{F}^{G}(U) = \mathbf{W}(\mathcal{F})^{G}(U)
= \bigl\{ x \in \mathcal{F}(U) \bigm|
g \cdot x_{S'} = x_{S'}
\text{ for every } S' \xrightarrow{f} U,\ g \in G(S') \bigr\},
\]
where \(x_{S'}\) denotes the image of \(x\) in
\(\Gamma(S', f^{*}(\mathcal{F})) = \Gamma(U, f_{*}f^{*}(\mathcal{F}))\).

One must take care that the natural morphism
\(\mathbf{W}(\mathcal{F}^{G}) \to \mathbf{W}(\mathcal{F})^{G}\) is not an isomorphism
in general. For example, if \(S = \Spec(\mathbb{Z})\) and \(G\) is the constant group
\(\mathbb{Z}/2\mathbb{Z} = \{1, \tau\}\) acting on \(\mathcal{F} = \OS\) by
\(\tau \cdot 1 = -1\), one has \(\mathcal{F}^{G} = 0\) but, if \(R\) is an
\(\mathbb{F}_{2}\)-algebra,
\(\mathbf{W}(\mathcal{F})^{G}(\Spec(R)) = R\).
\end{remark}

\numpara{4.7.2}
One assumes henceforth, until the end of \No4.7, that \(G\) is \emph{affine}
over~\(S\).%
{\renewcommand{\thefootnote}{(40)}\nde{Cf.\ \(\mathrm{VI}_{\mathrm{B}}\), §§~11.1--11.6 for the
extension of the results of 4.7.2 to the case where \(G\) is not necessarily
affine, but where \(G\) and \(\mathcal{F}\) are assumed \emph{flat} over~\(S\).}}
Then, by virtue of 4.6.4, the datum of a morphism of \(S\)-functors
\[
\rho \colon h_{G} \longrightarrow
\underline{\End}_{\mathbf{O}_{S}}\bigl(\mathbf{W}(\mathcal{F})\bigr)
\]
is equivalent to that of a morphism of \(\OS\)-modules
\[
\mu \colon \mathcal{F} \longrightarrow
\mathcal{F} \otimes_{\OS} \mathcal{A}(G).
\]

To say that \(\rho\) is a morphism of
\((\widehat{\mathbf{Sch}})_{/S}\)-groups is then equivalent to saying that
\(\mu\) satisfies the following axioms:

\medskip
\noindent
(CM~1) \emph{the following diagram is commutative}
\[
\xymatrix{
\mathcal{F} \ar[rr]^{\mu} \ar[d]_{\mu}
  && \mathcal{F} \otimes_{\OS} \mathcal{A}(G)
     \ar[d]^{\mathrm{id}\otimes\Delta} \\
\mathcal{F} \otimes_{\OS} \mathcal{A}(G) \ar[rr]^{\mu\otimes\mathrm{id}}
  && \mathcal{F} \otimes_{\OS} \mathcal{A}(G)
     \otimes_{\OS} \mathcal{A}(G)\,.
}
\]

\noindent
(CM~2) \emph{the composite below is the identity}
\[
\mathcal{F} \xrightarrow{\mu}
\mathcal{F} \otimes_{\OS} \mathcal{A}(G)
\xrightarrow{\mathrm{id}\otimes\varepsilon}
\mathcal{F} \otimes_{\OS} \OS
\xrightarrow{\sim}
\mathcal{F}.
\]

These axioms (CM~1) and (CM~2) are those of the structure of a (right)
comodule%
{\renewcommand{\thefootnote}{(41)}\nde{Left \(G\)-\(\OS\)-modules correspond naturally to right
\(\mathcal{A}(G)\)-comodules.}}
over the bialgebra \(\mathcal{A}(G)\).

Set \(\mathcal{A} = \mathcal{A}(G)\). If \(\mathcal{F}\) and \(\mathcal{F}'\)
are \(\mathcal{A}\)-comodules, a morphism of comodules
\(f \colon \mathcal{F} \to \mathcal{F}'\) is a morphism of \(\OS\)-modules such
that the following diagram is commutative:
\[
\xymatrix{
\mathcal{F} \ar[r]^{f} \ar[d]_{\mu_{\mathcal{F}}}
  & \mathcal{F}' \ar[d]^{\mu_{\mathcal{F}'}} \\
\mathcal{F} \otimes \mathcal{A} \ar[r]^{f\otimes\mathrm{id}}
  & \mathcal{F}' \otimes \mathcal{A}\,.
}
\]
One thus obtains the category \((\mathcal{A}\text{-}\mathrm{Comod.})\), and
one will denote by \((\mathcal{A}\text{-}\mathrm{Comod.q.c.})\) the full
subcategory consisting of the \(\mathcal{A}\)-comodules which are
quasi-coherent as \(\OS\)-modules.
One has thus obtained:

\begin{proposition}{4.7.2}\label{I.4.7.2}
--- Let \(G\) be an \(S\)-group affine over~\(S\). One has equivalences of
categories:
\begin{align*}
(G\text{-}\OS\text{-Mod.})
  &\simeq \bigl(\mathcal{A}(G)\text{-}\mathrm{Comod.}\bigr) \\
(G\text{-}\OS\text{-Mod.q.c.})
  &\simeq \bigl(\mathcal{A}(G)\text{-}\mathrm{Comod.q.c.}\bigr).
\end{align*}
{\renewcommand{\thefootnote}{(42)}\nde{The following sentence has been added.}}
If moreover \(S = \Spec(\Lambda)\) is affine and if one denotes
\(\Lambda[G] = \Gamma(S, \mathcal{A}(G))\), one has an equivalence of
categories
\[
\bigl(\mathcal{A}(G)\text{-}\mathrm{Comod.q.c.}\bigr)
\simeq \bigl(\Lambda[G]\text{-}\mathrm{Comod.}\bigr).
\]
\end{proposition}

{\renewcommand{\thefootnote}{(43)}\nde{The following paragraph has been added, taken from
\textup{[Se68, \S~1.3]}.}}
Suppose moreover that \(\mathcal{A} = \mathcal{A}(G)\) is a \emph{flat}
\(\OS\)-module. Let \(\mathcal{E}\) be an \(\mathcal{A}\)-comodule and
\(\mathcal{F}\) a sub-\(\OS\)-module of \(\mathcal{E}\). Since \(\mathcal{A}\)
is flat over \(\OS\), one can identify \(\mathcal{F} \otimes \mathcal{A}\)
(resp.\ \(\mathcal{F} \otimes \mathcal{A} \otimes \mathcal{A}\)) with a
sub-\(\OS\)-module of \(\mathcal{E} \otimes \mathcal{A}\)
(resp.\ \(\mathcal{E} \otimes \mathcal{A} \otimes \mathcal{A}\)).
Suppose that \(\mu_{\mathcal{E}}\) maps \(\mathcal{F}\) into
\(\mathcal{F} \otimes \mathcal{A}\); then its restriction
\(\mu_{\mathcal{F}} \colon \mathcal{F} \to \mathcal{F} \otimes \mathcal{A}\)
equips \(\mathcal{F}\) with a structure of \(\mathcal{A}\)-comodule; one says
that \(\mathcal{F}\) is a \emph{subcomodule} of \(\mathcal{E}\).
By passage to the quotient, \(\mu_{\mathcal{E}}\) defines a morphism of
\(\OS\)-modules
\(\mathcal{E}/\mathcal{F} \to \mathcal{E}/\mathcal{F} \otimes \mathcal{A}\),
which equips \(\mathcal{E}/\mathcal{F}\) with a structure of
\(\mathcal{A}\)-comodule. If \(f \colon \mathcal{E} \to \mathcal{E}'\) is a
morphism of \(\mathcal{A}\)-comodules, \(\Ker f\) (resp.\ \(\Im f\)) is a
sub-\(\mathcal{A}\)-comodule of \(\mathcal{E}\) (resp.\ \(\mathcal{E}'\)), and
\(f\) induces an isomorphism of \(\mathcal{A}\)-comodules:
\(\mathcal{E}/\Ker f \isomto \Im f\). Moreover, if \(\mathcal{E}\) and
\(\mathcal{E}'\) are quasi-coherent \(\OS\)-modules, the same is true of
\(\Ker f\) and \(\Im f\). Consequently, \((\mathcal{A}\text{-}\mathrm{Comod.})\)
and \((\mathcal{A}\text{-}\mathrm{Comod.q.c.})\) are \emph{abelian}
categories.

\begin{corollary}{4.7.2.1}\label{I.4.7.2.1}
--- One assumes that \(G\) is affine and flat over~\(S\). Then the category
\((G\text{-}\OS\text{-Mod.q.c.})\)
(resp.\ \((G\text{-}\OS\text{-Mod.})\)), equivalent to the
category of quasi-coherent \(\mathcal{A}(G)\)-comodules over \(\OS\)
(resp.\ of \(\mathcal{A}(G)\)-comodules), is abelian.
\end{corollary}

\numpara{4.7.3}
Suppose now that \(G\) is a \emph{diagonalizable} group, that is, that
\(\mathcal{A}(G)\) is the algebra of a commutative group \(M\) over the sheaf
of rings \(\OS\). If \(\mathcal{F}\) is an \(\OS\)-module, one has
\[
\mathcal{F} \otimes \mathcal{A}(G)
= \coprod_{m \in M} \mathcal{F} \otimes m\OS.
\]

To give a morphism of \(\OS\)-modules
\[
\mu \colon \mathcal{F} \longrightarrow \mathcal{F} \otimes \mathcal{A}(G)
\]
is therefore equivalent to giving \(\OS\)-endomorphisms
\((\mu_{m})_{m \in M}\) of \(\mathcal{F}\), such that for every section \(x\)
of \(\mathcal{F}\) over an open of \(S\), \((\mu_{m}(x))\) is a section of the
direct sum \(\coprod_{m \in M} \mathcal{F}\) (this means that over every
sufficiently small open, only finitely many of the restrictions of the
\(\mu_{m}(x)\) are nonzero).

In order that \(\mu\), defined by
\[
\mu(x) = \sum_{m \in M} \mu_{m}(x) \otimes m,
\]
satisfy (CM~1) (resp.\ (CM~2)), it is necessary and sufficient that one have
\[
\mu_{m} \circ \mu_{m'} = \delta_{mm'}\,\mu_{m},
\qquad
\Bigl(\text{resp.}\ \sum_{m \in M} \mu_{m} = \mathrm{Id}_{\mathcal{F}}\Bigr),
\]
which means that the \(\mu_{m}\) are pairwise orthogonal projectors summing to
the identity. One has thus proved:

\begin{proposition}{4.7.3}\label{I.4.7.3}
--- If \(G = \mathbf{D}_{S}(M)\) is a diagonalizable \(S\)-group, the category of
quasi-coherent \(G\)-\(\OS\)-modules
(resp.\ of \(G\)-\(\OS\)-modules) is equivalent to the category of
quasi-coherent \(\OS\)-modules (resp.\ of \(\OS\)-modules) graded of type~\(M\).
\end{proposition}

\begin{remark*}
--- If \(\mathcal{F}\) is endowed with a structure of \(\OS\)-algebra preserved by
the operations of \(G\), then the grading of \(\mathcal{F}\) is an algebra
grading. More precisely:
\end{remark*}

\begin{corollary}{4.7.3.1}\label{I.4.7.3.1}
--- The functor \(\mathcal{A} \mapsto \Spec \mathcal{A}\) induces an equivalence
between the category of quasi-coherent \(\OS\)-algebras graded of type~\(M\)
and the opposite category to that of the \(S\)-schemes affine over~\(S\) with
\(S\)-group of operators \(G = \mathbf{D}_{S}(M)\).
\end{corollary}

\begin{proposition}{4.7.4}\label{I.4.7.4}
--- Let \(G\) be a diagonalizable \(S\)-group. If
\(0 \to \mathcal{F}_{1} \to \mathcal{F}_{2} \to \mathcal{F}_{3} \to 0\)
is an exact sequence of quasi-coherent \(G\)-\(\OS\)-modules which
splits as a sequence of \(\OS\)-modules, then it likewise splits as a sequence
of \(G\)-\(\OS\)-modules.
\end{proposition}

Indeed, if \(G = \mathbf{D}_{S}(M)\), each of the \(\mathcal{F}_{i}\) is graded by the
\((\mathcal{F}_{i})_{m}\) and for each \(m \in M\) the sequence
\[
0 \longrightarrow (\mathcal{F}_{1})_{m} \longrightarrow
(\mathcal{F}_{2})_{m} \longrightarrow (\mathcal{F}_{3})_{m} \longrightarrow 0
\]
of \(\OS\)-modules is split. The preceding proposition then yields the result.

\sgasection{5}{Cohomology of groups}
\label{I.sec.5}

\par\addvspace{0.55\baselineskip}%
\noindent\textbf{5.1. The standard complex.}\ ---\ \label{I.5.1}%
{\renewcommand{\thefootnote}{(44)}\nde{This complex is often called the ``Hochschild complex''; see for example
\S~II.3 of \textup{[DG70]}.}}%
Let \(\mathcal{C}\) be a category, \(\mathbf{G}\) a
\(\widehat{\mathcal{C}}\)-group, \(\mathbf{O}\) a
\(\widehat{\mathcal{C}}\)-ring and \(\mathbf{F}\) a
\(\mathbf{G}\)-\(\mathbf{O}\)-module. One sets, for \(n \geqslant 0\),
\[
\underline{C}^{n}(\mathbf{G}, \mathbf{F})
= \uHom(\mathbf{G}^{n}, \mathbf{F})
\quad\text{and}\quad
C^{n}(\mathbf{G}, \mathbf{F})
= \Hom(\mathbf{G}^{n}, \mathbf{F}),
\]
where \(\mathbf{G}^{0}\) is the final object~\(e\). Then
\(\underline{C}^{n}(\mathbf{G}, \mathbf{F})\)
(resp.\ \(C^{n}(\mathbf{G}, \mathbf{F})\)) is endowed in an evident way with a
structure of \(\mathbf{O}\)-module
(resp.\ of \(\Gamma(\mathbf{O})\)-module), and one has
\[
C^{n}(\mathbf{G}, \mathbf{F})
\simeq \Gamma\bigl(\underline{C}^{n}(\mathbf{G}, \mathbf{F})\bigr)
\quad\text{and}\quad
\underline{C}^{n}(\mathbf{G}, \mathbf{F})(S)
= C^{n}(\mathbf{G}_{S}, \mathbf{F}_{S}).
\]

To give an element of \(C^{n}(\mathbf{G}, \mathbf{F})\) is to give, for each
\(S \in \Ob(\mathcal{C})\), an \(n\)-cochain of \(\mathbf{G}(S)\) in
\(\mathbf{F}(S)\), functorially in~\(S\). The boundary operator
\[
\partial \colon C^{n}\bigl(\mathbf{G}(S), \mathbf{F}(S)\bigr)
\longrightarrow C^{n+1}\bigl(\mathbf{G}(S), \mathbf{F}(S)\bigr),
\]
which, let us recall, is given by the formula
\begin{align*}
\partial f(g_{1}, \ldots, g_{n+1})
&= g_{1} f(g_{2}, \ldots, g_{n+1})
+ \sum_{i=1}^{n} (-1)^{i}
f(g_{1}, \ldots, g_{i}g_{i+1}, \ldots, g_{n+1}) \\
&\qquad + (-1)^{n+1} f(g_{1}, \ldots, g_{n}),
\end{align*}
is functorial in \(S\) and therefore defines a homomorphism
\[
\partial \colon C^{n}(\mathbf{G}, \mathbf{F})
\longrightarrow C^{n+1}(\mathbf{G}, \mathbf{F})
\]
such that \(\partial \circ \partial = 0\). One has thus defined a complex of
abelian groups (and even of \(\Gamma(\mathbf{O})\)-modules), denoted
\(C^{*}(\mathbf{G}, \mathbf{F})\). One defines in the same way the complex of
\(\mathbf{O}\)-modules \(\underline{C}^{*}(\mathbf{G}, \mathbf{F})\), and one
has
\[
C^{*}(\mathbf{G}, \mathbf{F})
= \Gamma\bigl(\underline{C}^{*}(\mathbf{G}, \mathbf{F})\bigr).
\]
One denotes by \(H^{n}(\mathbf{G}, \mathbf{F})\)
(resp.\ \(\underline{H}^{n}(\mathbf{G}, \mathbf{F})\)) the homology groups
(resp.\ the homology \(\widehat{\mathcal{C}}\)-groups) of the complex
\(C^{*}(\mathbf{G}, \mathbf{F})\)
(resp.\ \(\underline{C}^{*}(\mathbf{G}, \mathbf{F})\)).

One has in particular
\[
\underline{H}^{0}(\mathbf{G}, \mathbf{F}) = \mathbf{F}^{\mathbf{G}}
\quad\text{and}\quad
H^{0}(\mathbf{G}, \mathbf{F}) = \Gamma(\mathbf{F}^{\mathbf{G}}).
\]

\begin{remark}{5.1.1}\label{I.5.1.1}
--- {\renewcommand{\thefootnote}{(45)}\nde{This remark has been added.}}
The ``set-theoretic'' description of \(\partial\) given above is convenient
for checking that \(\partial \circ \partial = 0\). One can also define
\(\partial\) in terms of the multiplication
\(m \colon \mathbf{G} \times \mathbf{G} \to \mathbf{G}\) and of the action
\(\mu \colon \mathbf{G} \times \mathbf{F} \to \mathbf{F}\) as follows: for every
\(f \in C^{n}(\mathbf{G}, \mathbf{F})\),
\[
\partial f = \mu \circ (\mathrm{id}_{\mathbf{G}} \times f)
+ \sum_{i=1}^{n} (-1)^{i}
f \circ \bigl(\mathrm{id}_{\mathbf{G}^{i-1}} \times m
\times \mathrm{id}_{\mathbf{G}^{n-i}}\bigr)
+ (-1)^{n+1} f \circ \mathrm{pr}_{[1,n]},
\]
where \(\mathrm{pr}_{[1,n]}\) denotes the projection of
\(\mathbf{G}^{n+1} = \mathbf{G}^{n} \times \mathbf{G}\) onto
\(\mathbf{G}^{n}\). Likewise, for every \(S \in \Ob(\mathcal{C})\) and
\(f \in \underline{C}^{n}(\mathbf{G}, \mathbf{F})(S)
= C^{n}(\mathbf{G}_{S}, \mathbf{F}_{S})\), one has
\[
\partial f = \mu_{S} \circ (\mathrm{id}_{\mathbf{G}_{S}} \times f)
+ \sum_{i=1}^{n} (-1)^{i}
f \circ \bigl(\mathrm{id}_{\mathbf{G}_{S}^{i-1}} \times m_{S}
\times \mathrm{id}_{\mathbf{G}_{S}^{n-i}}\bigr)
+ (-1)^{n+1} f \circ \mathrm{pr}_{[1,n]},
\]
where \(m_{S}\) and \(\mu_{S}\) are deduced from \(m\) and \(\mu\) by base
change.
\end{remark}

\numpara{5.2}\label{I.5.2}
{\renewcommand{\thefootnote}{(46)}\nde{The following reminder has been added.}}
One recalls (\cf \S~3) that \((\mathbf{G}\text{-}\mathbf{O}\text{-}\mathrm{Mod.})\)
is endowed with the structure of an abelian category, defined ``argument by
argument''; thus,
\[
0 \longrightarrow \mathbf{F}' \longrightarrow \mathbf{F}
\longrightarrow \mathbf{F}'' \longrightarrow 0
\]
is an exact sequence of \(\mathbf{G}\)-\(\mathbf{O}\)-modules if, and only if,
the sequence of abelian groups
\[
0 \longrightarrow \mathbf{F}'(S) \longrightarrow \mathbf{F}(S)
\longrightarrow \mathbf{F}''(S) \longrightarrow 0
\]
is exact, for every \(S \in \Ob(\mathcal{C})\).

{\renewcommand{\thefootnote}{(47)}\nde{The following sentence has been added.}}
Suppose \(\mathcal{C}\) is small; then, by 3.2.1,
\((\mathbf{G}\text{-}\mathbf{O}\text{-}\mathrm{Mod.})\) has enough injective
objects, so that the derived functors of the left exact functors \(H^{0}\) and
\(\underline{H}^{0}\) are defined. We now propose to show that the functors
\(H^{n}\) (resp.\ \(\underline{H}^{n}\)) are indeed the derived functors of
\(H^{0}\) (resp.\ \(\underline{H}^{0}\)).

\begin{definition}{5.2.0}\label{I.5.2.0}
--- {\renewcommand{\thefootnote}{(48)}\nde{The original has been modified, in order to introduce 5.2.0.1 and
5.2.0.2, which will be useful in the proof of Theorem~5.3.1.}}%
For every \(\mathbf{O}\)-module \(\mathbf{P}\), one denotes by \(E(\mathbf{P})\)
the object \(\uHom(\mathbf{G}, \mathbf{P})\) of \(\widehat{\mathcal{C}}\)
endowed with the structure of \(\mathbf{G}\)-\(\mathbf{O}\)-module defined as
follows: for every \(S \in \Ob(\mathcal{C})\) one has
\(\uHom(\mathbf{G}, \mathbf{P})(S)
= \Hom_{S}(\mathbf{G}_{S}, \mathbf{P}_{S})\), and one lets
\(g \in \mathbf{G}(S)\) and \(a \in \mathbf{O}(S)\) operate on
\(\phi \in \Hom_{S}(\mathbf{G}_{S}, \mathbf{P}_{S})\) by the formulas
\[
(g \cdot \phi)(h) = \phi(hg)
\quad\text{and}\quad
(a \cdot \phi)(h) = a\phi(h),
\]
for every \(h \in \mathbf{G}(S')\), \(S' \to S\). Moreover, for every
\(\phi \in \Hom_{S}(\mathbf{G}_{S}, \mathbf{P}_{S})\) one sets
\[
\varepsilon(\phi) = \phi(1) \in \mathbf{P}(S)
\]
(where \(1\) denotes the unit element of \(\mathbf{G}(S)\)).

This defines a functor
\(E \colon (\mathbf{O}\text{-}\mathrm{Mod.}) \to
(\mathbf{G}\text{-}\mathbf{O}\text{-}\mathrm{Mod.})\)
and a natural transformation \(\varepsilon \colon E \to \mathrm{Id}\), where
\(\mathrm{Id}\) denotes the identity functor of
\((\mathbf{O}\text{-}\mathrm{Mod.})\).
\end{definition}

\begin{remark}{5.2.0.1}\label{I.5.2.0.1}
--- {\renewcommand{\thefootnote}{(48)}\footnotemark[\value{footnote}]}
In what follows, let us denote by \(\mathbf{G}_{1}\) and \(\mathbf{G}_{2}\)
two copies of \(\mathbf{G}\). Then the morphism
\[
\mathbf{G}_{1} \times E(\mathbf{P}) \longrightarrow E(\mathbf{P}),
\qquad
(g_{1}, \phi) \longmapsto \bigl(g_{2} \mapsto \phi(g_{2}g_{1})\bigr)
\]
corresponds, via the isomorphisms
\begin{align*}
\Hom\bigl(\mathbf{G}_{1} \times E(\mathbf{P}), E(\mathbf{P})\bigr)
&\simeq \Hom\bigl(E(\mathbf{P}),
\uHom\bigl(\mathbf{G}_{1}, \uHom(\mathbf{G}_{2}, \mathbf{P})\bigr)\bigr) \\
&\simeq \Hom\bigl(E(\mathbf{P}),
\uHom(\mathbf{G}_{2} \times \mathbf{G}_{1}, \mathbf{P})\bigr)
\end{align*}
to the morphism
\(\phi \mapsto \bigl((g_{2}, g_{1}) \mapsto \phi(g_{2}g_{1})\bigr)\),
\ie to the morphism
\[
\uHom(\mathbf{G}, \mathbf{P}) \longrightarrow
\uHom(\mathbf{G}_{2} \times \mathbf{G}_{1}, \mathbf{P})
\]
induced by the multiplication
\(\mu_{\mathbf{G}} \colon \mathbf{G} \times \mathbf{G} \to \mathbf{G}\),
\((g_{2}, g_{1}) \mapsto g_{2}g_{1}\).
\end{remark}

\begin{lemma}{5.2.0.2}\label{I.5.2.0.2}
--- {\renewcommand{\thefootnote}{(48)}\footnotemark[\value{footnote}]}
\par
\begin{enumeratei}
\item The functor \(E\) is right adjoint to the forgetful functor
\((\mathbf{G}\text{-}\mathbf{O}\text{-}\mathrm{Mod.}) \to
(\mathbf{O}\text{-}\mathrm{Mod.})\); more precisely,
\(\varepsilon \colon E \to \mathrm{Id}\) induces, for every
\[
\mathbf{M} \in (\mathbf{G}\text{-}\mathbf{O}\text{-}\mathrm{Mod.})
\quad\text{and}\quad
\mathbf{P} \in \Ob(\mathbf{O}\text{-}\mathrm{Mod.}),
\]
a bijection
\[
\Hom_{\mathbf{G}\text{-}\mathbf{O}\text{-}\mathrm{Mod.}}
\bigl(\mathbf{M}, E(\mathbf{P})\bigr)
\xrightarrow{\sim}
\Hom_{\mathbf{O}\text{-}\mathrm{Mod.}}\bigl(\mathbf{M}, \mathbf{P}\bigr)
\]
functorial in \(\mathbf{M}\) and \(\mathbf{P}\).
\item Consequently, if \(\mathbf{I}\) is an injective object of
\((\mathbf{O}\text{-}\mathrm{Mod.})\) then \(E(\mathbf{I})\) is an injective
object of \((\mathbf{G}\text{-}\mathbf{O}\text{-}\mathrm{Mod.})\).
\end{enumeratei}
\end{lemma}

\begin{proof}
To every \(\mathbf{O}\)-morphism \(f \colon \mathbf{M} \to \mathbf{P}\), one
associates the element \(\phi_{f}\) of
\[
\Hom_{\mathbf{O}}\bigl(\mathbf{M}, E(\mathbf{P})\bigr)
\]
defined as follows.
For every \(S \in \Ob(\mathcal{C})\) and \(m \in \mathbf{M}(S)\),
\(\phi_{f}(m)\) is the element of
\(\Hom_{S}(\mathbf{G}_{S}, \mathbf{P}_{S})\) defined by: for every
\(g \in \mathbf{G}(S')\), \(S' \to S\),
\[
\phi_{f}(m)(g) = f(gm) \in \mathbf{P}(S').
\]
Then, for every \(h \in \mathbf{G}(S)\), one has
% typo?: kept as printed; expected h\cdot\varphi_f(m)
\(\phi_{f}(hm) = h \cdot f(m)\),
\ie \(\phi_{f}\) is an element of
\[
\Hom_{\mathbf{G}\text{-}\mathbf{O}\text{-}\mathrm{Mod.}}
\bigl(\mathbf{M}, E(\mathbf{P})\bigr).
\]
If
\[
\phi \in \Hom_{\mathbf{G}\text{-}\mathbf{O}\text{-}\mathrm{Mod.}}
\bigl(\mathbf{M}, E(\mathbf{P})\bigr)
\]
and if one denotes, for every
\(m \in \mathbf{M}(S)\), \(f(m) = \phi(m)(1)\), then
\[
\phi_{f}(m)(g) = f(gm) = \phi(gm)(1)
= \bigl(g \cdot \phi(m)\bigr)(1) = \phi(m)(g),
\]
\ie \(\phi_{f} = \phi\). Conversely, it is clear that
\(\phi_{f}(m)(1) = f(m)\). This proves~(i), and~(ii) follows from it at once.
\end{proof}

\begin{definition}{5.2.0.3}\label{I.5.2.0.3}
--- Let \(\mathbf{M}\) be a \(\mathbf{G}\)-\(\mathbf{O}\)-module; the identity map
of \(\mathbf{M}\) (considered as an \(\mathbf{O}\)-module) corresponds by
adjunction to the morphism of \(\mathbf{G}\)-\(\mathbf{O}\)-modules
\[
\mu_{\mathbf{M}} \colon \mathbf{M} \longrightarrow E(\mathbf{M})
\]
such that for every \(S \in \Ob(\mathcal{C})\) and \(m \in \mathbf{M}(S)\),
\(\mu_{\mathbf{M}}(m)\) is the morphism
\(\mathbf{G}_{S} \to \mathbf{M}_{S}\) defined by: for every \(S' \to S\) and
\(g \in \mathbf{G}(S')\),
\(\mu_{\mathbf{M}}(m)(g) = g \cdot m_{S'} \in \mathbf{M}(S')\).

Note that \(\mu_{\mathbf{M}}\) is a \emph{monomorphism}. Indeed,
\(\varepsilon_{\mathbf{M}} \colon E(\mathbf{M}) \to \mathbf{M}\) is a morphism
of \(\mathbf{O}\)-modules such that
\(\varepsilon_{\mathbf{M}} \circ \mu_{\mathbf{M}} = \mathrm{id}_{\mathbf{M}}\);
consequently \(\mathbf{M}\) is a direct factor, as an \(\mathbf{O}\)-module,
of \(E(\mathbf{M})\).
\end{definition}

\begin{proposition}{5.2.1}\label{I.5.2.1}
--- One assumes that \(\mathcal{C}\) is small, that finite products exist in it,
and that \(\mathbf{G}\) is representable. Then the functors
\(H^{n}(\mathbf{G},\,)\) (resp.\ \(\underline{H}^{n}(\mathbf{G},\,)\)) are the
derived functors of the left exact functor \(H^{0}(\mathbf{G},\,)\)
(resp.\ \(\underline{H}^{0}(\mathbf{G},\,)\)) on the category of
\(\mathbf{G}\)-\(\mathbf{O}\)-modules.
\end{proposition}

By virtue of the well-known general results%
{\renewcommand{\thefootnote}{(49)}\nde{Cf.\ \textup{[Gr57]}, 2.2.1
and 2.3. Moreover, the original has been spelled out in what follows.}},
it suffices to check that the \(H^{n}(\mathbf{G},\,)\)
(resp.\ \(\underline{H}^{n}(\mathbf{G},\,)\)) form an effaceable cohomological
functor in degrees \(> 0\).

Let
\[
0 \longrightarrow \mathbf{F}' \longrightarrow \mathbf{F}
\longrightarrow \mathbf{F}'' \longrightarrow 0
\]
be an exact sequence of \(\mathbf{G}\)-\(\mathbf{O}\)-modules, and let
\(S \in \Ob(\mathcal{C})\). By hypothesis, \(\mathbf{G}\) is representable by
an object \(G \in \Ob(\mathcal{C})\), and finite products exist in
\(\mathcal{C}\); in particular \(\mathcal{C}\) has a final element~\(e\).
Thus each \(G^{n} \times \mathbf{h}_{S}\) is representable by \(G^{n} \times S\)
(with \(G^{0} = e\)), and the sequence
\[
0 \longrightarrow \mathbf{F}'(G^{n} \times S) \longrightarrow
\mathbf{F}(G^{n} \times S) \longrightarrow
\mathbf{F}''(G^{n} \times S) \longrightarrow 0
\]
is exact. This shows that the sequence of \(\mathbf{O}\)-modules
\[
0 \longrightarrow \underline{C}^{n}(\mathbf{h}_{G}, \mathbf{F}')
\longrightarrow \underline{C}^{n}(\mathbf{h}_{G}, \mathbf{F})
\longrightarrow \underline{C}^{n}(\mathbf{h}_{G}, \mathbf{F}'')
\longrightarrow 0
\]
is exact. It follows that \(\underline{C}^{*}(\mathbf{G},\,)\), considered as
a functor on \((\mathbf{G}\text{-}\mathbf{O}\text{-}\mathrm{Mod.})\) with
values in the category of complexes of
\((\mathbf{O}\text{-}\mathrm{Mod.})\), is exact. This shows that the
\(\underline{H}^{n}(\mathbf{G},\,)\) do indeed form a cohomological functor.
